\section{Fundamentação Teórica}

\subsection{Segurança da Informação e Ameaças em Rede}
A segurança da informação tornou-se elemento central na sociedade contemporânea
devido à crescente dependência de sistemas digitais. Com a consolidação da chamada
“civilização cibernética” \cite{wasyhun2023cyber}, iniciada com a revolução
digital no final do século XX \cite{Kazi2019network}, houve aumento significativo
na produção e circulação de dados, ampliando a superfície de ataque. O foco da área
é proteger dados contra ameaças internas e externas, assegurando a tríade CIA de
confidencialidade (\en{confidentiality}), integridade (\en{integrity}) e
disponibilidade (\en{availability}) \cite{diptiben2022cyber}. Ataques
cibernéticos podem ocasionar perdas financeiras, danos à reputação e impactos
críticos quando afetam infraestruturas essenciais \cite{wasyhun2023cyber}.

As ameaças em redes são classificadas conforme objetivo e forma de atuação. Entre
as principais estão a negação de serviço (DoS), que visa indisponibilizar
recursos; a sondagem (\en{probing}), que coleta informações da rede; o
escalonamento de privilégios (U2R) e o acesso remoto indevido (R2L)
\cite{wasyhun2023cyber}. Somam-se a elas \en{malware}, \en{botnets} e engenharia
social \cite{Thakkar2021Survey, solane2015using}. Um dos maiores desafios são os
ataques \en{zero-day}, que exploram vulnerabilidades ainda desconhecidas e, por não
possuírem assinatura, são difíceis de identificar por métodos tradicionais
\cite{wasyhun2023cyber}.

\subsection{Sistemas de Detecção de Intrusão}
Os IDS podem ser classificados pela forma de implantação e pelo método de detecção.
Quanto à implantação, os sistemas baseados em \en{host} (HIDS) monitoram atividades
locais e os baseados em rede (NIDS) analisam o tráfego em pontos estratégicos da
infraestrutura. Quanto ao método, os IDS operam por assinatura, identificando
padrões conhecidos, ou por anomalia, detectando desvios em relação ao comportamento
considerado normal; abordagens híbridas combinam as duas
\cite{Khraisat2019Survey, zehedi2023comparative, Thakkar2021Survey, Khraisat2019Hybrid}.
Sistemas por assinatura são incapazes de identificar ataques \en{zero-day},
enquanto sistemas por anomalia tendem a gerar altas taxas de falso positivo ao
interpretar comportamentos legítimos como ameaças \cite{Kazi2019network}.

Sommer e Paxson \cite{Sommer2010Outside} argumentam que a detecção de intrusão
difere de outros domínios em que o aprendizado de máquina tem sucesso: o custo de
um erro é muito alto, o tráfego “normal” é extremamente variável, faltam bases de
treino representativas e há uma lacuna semântica entre “anômalo” e “malicioso”.
Por isso recomendam que o sistema explique suas decisões ao operador e que a
avaliação vá além da acurácia em bases fechadas. Esses pontos orientaram
diretamente as decisões do MVP2.

\subsection{Aprendizado de Máquina em IDS}
O ML permite a análise automatizada de grandes volumes de dados de rede e a
identificação de padrões complexos sem intervenção humana constante
\cite{zehedi2023comparative, Thakkar2021Survey}. As técnicas dividem-se em
aprendizado supervisionado, não supervisionado e por reforço, e o processo envolve
a divisão dos dados em conjuntos de treinamento e teste para avaliação
\cite{zehedi2023comparative, emad2022intrusion}. Entre os desafios estão o custo
computacional, o desbalanceamento das classes --- que prejudica a detecção de
ataques raros --- e a dependência da qualidade dos dados \cite{zehedi2023comparative}.
Etapas de pré-processamento, seleção de atributos e redução de dimensionalidade,
como o PCA, são essenciais para o desempenho, e a avaliação usa métricas como
acurácia, precisão, \en{recall} e \en{F1-score} \cite{zehedi2023comparative, Thakkar2021Survey}.

\subsection{Trabalhos Relacionados}
Liao \textit{et al.} \cite{Liao2013Intrusion} revisam técnicas de IDS baseadas em
Bayes, teoria dos jogos e mineração de dados, e mostram que árvores de decisão e SVM
apresentam altas taxas de acurácia, sem uma abordagem dominante. Saranya
\textit{et al.} \cite{Saranya2020Performance} comparam algoritmos em cenários IoT e
5G e obtêm 99,65\% de acurácia com \en{Random Forest}, contra 98\% do CART isolado.
Khraisat \textit{et al.} \cite{Khraisat2019Hybrid} combinam C5.0 e
\en{One-Class SVM} em um modelo híbrido, com 83,24\% de acurácia no NSL-KDD e
97,04\% no ADFA, mostrando que a combinação entre detecção por assinatura e por
anomalia supera as limitações individuais --- ideia próxima da cascata adotada no
MVP2.

A interpretabilidade das árvores de decisão é destacada por Azam \textit{et al.}
\cite{zehedi2023comparative}, em que o CART alcançou 97,36\% de sucesso. Depren
\textit{et al.} \cite{Depren2005Intelligent} combinam C4.5 com mapas
auto-organizáveis (SOM) e atingem 99,84\% na detecção de uso indevido. Gao
\textit{et al.} \cite{Xianwei2019Adaptive} propõem o \en{MultiTree}, um
\en{ensemble} de árvores CART com votação majoritária que trata o desbalanceamento
e alcança 85,2\% em cenários desafiadores. Rathore \textit{et al.}
\cite{Rathore2016RealTime} usam C4.5 e REPTree em ambiente de \en{Big Data} com
apenas 9 atributos e obtêm taxa de verdadeiros positivos acima de 99,9\%.
Abdulganiyu \textit{et al.} \cite{Abdulganiyu2023Systematic} consolidam que
algoritmos baseados em árvores variam entre 85,2\% e 99,9\% de acurácia, conforme a
preparação dos dados.

Quanto ao pré-processamento, Vanin \textit{et al.} \cite{patrick2022study}
destacam o alto número de atributos do KDDCup99 (41) e do CICIDS2017 (80) e a
necessidade de técnicas como o PCA. Thakkar e Lohiya \cite{Thakkar2021Survey}
mostram que reduzir para 5 a 20 atributos relevantes melhora a eficiência dos
algoritmos. Iglesias e Zseby \cite{felix2014analysis} reduzem de 41 para 16
atributos sem perda de acurácia, e Othman \textit{et al.} \cite{Othman2018Intrusion}
selecionam 17 atributos por qui-quadrado e obtêm AUROC de 99,55\%. Abdallah
\textit{et al.} \cite{emad2022intrusion} reforçam que conjuntos de 10 a 19
atributos são comuns em abordagens eficientes. A Tabela~\ref{tab:revisao}, nos
Anexos, organiza os trabalhos estudados.

\subsection{Isolation Forest}
O \en{Isolation Forest} \cite{Liu2008Isolation} detecta anomalias isolando
pontos em vez de modelar a densidade do normal. Cada árvore particiona os dados
com cortes aleatórios até isolar cada amostra; pontos raros e distantes são
isolados com poucos cortes. O escore de anomalia de uma amostra $x$ em um conjunto
de $n$ amostras é

\begin{equation}
  s(x,n) = 2^{-\frac{E[h(x)]}{c(n)}},
  \label{eq:if}
\end{equation}

\noindent em que $E[h(x)]$ é a profundidade média de isolamento de $x$ nas árvores
e $c(n)$ é a profundidade média de uma busca malsucedida em árvore binária, usada
como normalização. Escores próximos de 1 indicam anomalia. O parâmetro
\en{contamination} define a fração de amostras de treino marcadas como anômalas e,
portanto, o limiar de decisão. Uma consequência estrutural, explorada na Seção IV,
é que um conjunto \emph{volumoso e homogêneo} de amostras de ataque deixa de ser
raro e passa a ser difícil de isolar.

\subsection{Máquina de Vetores de Suporte}
A SVM \cite{Cortes1995Support} procura o hiperplano $f(x)=w^\top x + b$ de
máxima margem entre duas classes. Na versão linear, o valor de $f(x)$ é uma
distância com sinal ao hiperplano (a menos do fator $\lVert w\rVert$): o sinal
dá a classe predita e o módulo dá a confiança. Essa propriedade é usada no
Estágio 2 para ordenar a fila de revisão do analista. Por ser supervisionada, a
SVM só reconhece formas de ataque que já viu rotuladas.

\subsection{Explicabilidade por Modelo Substituto}
Um modelo substituto (\en{surrogate}) é um modelo interpretável treinado para
imitar as saídas de um modelo opaco \cite{Craven1996Extracting}. Treinando uma
árvore de decisão rasa sobre os rótulos atribuídos pelo \en{Isolation Forest},
obtêm-se regras legíveis, como \cod{n\_dports > 50 e duration <= 2,1}, que
acompanham cada alerta. A fidelidade do substituto mede o quanto ele concorda com
o modelo original --- e não o quanto o modelo original acerta.

\subsection{Envenenamento de Modelos e Aprendizado Ativo}
Ataques de envenenamento (\en{poisoning}) manipulam os dados de treino para
degradar ou desviar um modelo \cite{Biggio2012Poisoning}. Em um IDS que se
retreina sobre o tráfego observado, o atacante não precisa injetar dados
artificialmente: basta que sua própria atividade, mantida em volume baixo e
regular, entre no conjunto de treino e passe a ser aprendida como normal.

O aprendizado ativo (\en{active learning}) \cite{Settles2009Active} reduz o custo
de rotulagem escolhendo quais amostras submeter ao especialista. A estratégia por
incerteza prioriza as amostras mais próximas da fronteira de decisão, nas quais um
rótulo humano mais altera o modelo.
